\documentclass{article}

\usepackage{amsmath,amssymb}
\usepackage{xurl}
\usepackage[hidelinks]{hyperref}
\setlength{\emergencystretch}{2em}

% Shared commands for notes in docs/paper-gaps/.

\DeclareUnicodeCharacter{2080}{\ifmmode{}_{0}\else\textsubscript{0}\fi}
\DeclareUnicodeCharacter{2081}{\ifmmode{}_{1}\else\textsubscript{1}\fi}
\DeclareUnicodeCharacter{2082}{\ifmmode{}_{2}\else\textsubscript{2}\fi}
\DeclareUnicodeCharacter{2099}{\ifmmode{}_{n}\else\textsubscript{n}\fi}

\newcommand{\Lean}{\textsc{Lean}}
\ExplSyntaxOn
\NewDocumentCommand{\leanid}{m}{
  \ifmmode
    \textsf{\detokenize{#1}}
  \else
    \group_begin:
    \urlstyle{sf}
    \tl_set:Nn \l_tmpa_tl {#1}
    \tl_replace_all:Nnn \l_tmpa_tl {\_} {_}
    \exp_args:NV \path \l_tmpa_tl
    \group_end:
  \fi
}
\ExplSyntaxOff
\providecommand{\tr}{\operatorname{tr}}
\newcommand{\tnleanIssueBase}{https://github.com/LionSR/TNLean/issues/}
\newcommand{\tnleanPrBase}{https://github.com/LionSR/TNLean/pull/}
\newcommand{\tnleanDocBase}{https://sirui-lu.com/TNLean/docs/}
\newcommand{\ghissue}[1]{\href{\tnleanIssueBase#1}{GitHub issue \##1}}
\newcommand{\ghpr}[1]{\href{\tnleanPrBase#1}{PR \##1}}
\newcommand{\doclink}[1]{\href{\tnleanDocBase#1.html}{\path{#1}}}

% Dirac notation and the identity operator, as in blueprint/src/macros/common.tex.
% \providecommand keeps note-local definitions authoritative.
\usepackage{dsfont}
\providecommand{\ket}[1]{|#1\rangle}
\providecommand{\bra}[1]{\langle #1|}
\providecommand{\braket}[2]{\langle #1 | #2 \rangle}
\providecommand{\ketbra}[2]{|#1\rangle\!\langle #2|}
\providecommand{\Id}{\mathds{1}}
\providecommand{\one}{\mathds{1}}
\providecommand{\C}{\mathbb{C}}
\providecommand{\MN}[1]{M_{#1}(\C)}
\providecommand{\MD}{M_D(\C)}

% External referencing for paper sources.  The build helper writes sanitized
% auxiliary files under build/paper-aux/ and excludes duplicated source labels.
\usepackage{xr}

% Import a generated paper auxiliary file with a caller-supplied label prefix.
% The candidate paths support builds from docs/paper-gaps/, the repository root,
% and build/paper-aux/ itself.
\newcommand{\importpaperaux}[2]{%
  \IfFileExists{../../build/paper-aux/#2.aux}{%
    \externaldocument[#1]{../../build/paper-aux/#2}%
  }{%
    \IfFileExists{build/paper-aux/#2.aux}{%
      \externaldocument[#1]{build/paper-aux/#2}%
    }{%
      \IfFileExists{#2.aux}{%
        \externaldocument[#1]{#2}%
      }{}%
    }%
  }%
}

% General external reference macros
\newcommand{\paperref}[2]{\ref{#1-#2}}
\newcommand{\papereqref}[2]{(\ref{#1-#2})}

% CPSV16 source (1606.00608 / MPDO-22-12-17-2.tex) external document and shortcuts
\importpaperaux{cpsv16-}{1606.00608/MPDO-22-12-17-2-tnlean}
\newcommand{\cpsvref}[1]{\paperref{cpsv16}{#1}}
\newcommand{\cpsveqref}[1]{\papereqref{cpsv16}{#1}}

% Verdict marker: \gapnote{<kind>}{<status>}. Kind is the policy
% classification (clarification, local-correction, scope-restriction,
% unfaithful, false-source, open-gap); status is open, wip (elimination
% underway), resolved, or historical. Machine-read by the site index;
% prints a verdict line.
\newcommand{\gapnote}[2]{\par\noindent\textbf{Verdict:} #1 (#2).\par}


\title{Finite-Ring Entropy Corrections for CPSV16 Examples 4.10 and 4.11}
\author{}
\date{2026-08-12}

\begin{document}
\maketitle
\gapnote{false-source}{resolved}

\section{Scope}

Examples~4.10 and~4.11 of Cirac--P\'erez-Garc\'ia--Schuch--Verstraete appear
at lines~897--924 of
\path{Papers/1606.00608/MPDO-22-12-17-2.tex}
\cite{Cirac2016MPDO_arXiv}. This note distinguishes the literal source
statements from two finite-ring corrections.

The calculations below use base-two logarithms, matching the numerical values
printed in Example~4.10 of Ref.~\cite{Cirac2016MPDO_arXiv}. The formal entropy
definition uses natural logarithms. Dividing an entropy or mutual information
in nats by $\ln 2$ converts it to bits.

For a normalized four-spin state, let $S_L$ denote the entropy of $L$
consecutive spins. Define
\begin{align}
    I_L
        &= S_L+S_{4-L}-S_4,
    \label{eq:cpsv16_examples_4_10_4_11_entropy_mutual_information}\\
    I_2-I_1
        &= 2S_2-S_1-S_3.
    \label{eq:cpsv16_examples_4_10_4_11_entropy_mutual_information_defect}
\end{align}

\section{Example~4.10: the correlated bit flip}

\subsection{Printed channel and intended correction}

Each physical spin contains a left qubit and a right qubit. The source prints
the conjugating operator with the left-qubit label in both tensor factors:
\begin{align}
    \sigma_x^{(n,l)}\otimes\sigma_x^{(n,l)}.
    \label{eq:cpsv16_examples_4_10_4_11_entropy_printed_flip}
\end{align}
This expression does not act once on each qubit of spin~$n$. The
network-consistent correction is
\begin{align}
    U_n
        &= \sigma_x^{(n,l)}\otimes\sigma_x^{(n,r)},
    \label{eq:cpsv16_examples_4_10_4_11_entropy_corrected_flip}\\
    \mathcal E_n(X)
        &= pU_nXU_n^\dagger+(1-p)X.
    \label{eq:cpsv16_examples_4_10_4_11_entropy_corrected_channel}
\end{align}
At $p=1/4$, this correction reproduces all four entropy values printed after
the channel in Example~4.10 of Ref.~\cite{Cirac2016MPDO_arXiv}. The bundled
source states no alternative corrected channel.

\subsection{Exact four-spin calculation at $p=1/4$}

Start with a periodic ring of four Bell pairs, one on each bond, and apply
Eq.~\ref{eq:cpsv16_examples_4_10_4_11_entropy_corrected_channel}
independently at every spin. The reduced density operators for one through four
consecutive spins have the following nonzero spectra, where $a^{[m]}$ denotes
an eigenvalue $a$ with multiplicity $m$:
\begin{align}
    \operatorname{spec}(\rho_1)
        &= \left\{(1/4)^{[4]}\right\},
    \label{eq:cpsv16_examples_4_10_4_11_entropy_example410_spectrum_one}\\
    \operatorname{spec}(\rho_2)
        &= \left\{(5/32)^{[4]},(3/32)^{[4]}\right\},
    \label{eq:cpsv16_examples_4_10_4_11_entropy_example410_spectrum_two}\\
    \operatorname{spec}(\rho_3)
        &= \left\{(7/64)^{[4]},(3/64)^{[12]}\right\},
    \label{eq:cpsv16_examples_4_10_4_11_entropy_example410_spectrum_three}\\
    \operatorname{spec}(\rho_4)
        &= \left\{(41/128)^{[1]},(15/128)^{[4]},
            (9/128)^{[3]}\right\}.
    \label{eq:cpsv16_examples_4_10_4_11_entropy_example410_spectrum_four}
\end{align}
Each spectrum sums to one. Substitution into
$S(\rho)=-\sum_\lambda\lambda\log_2\lambda$ gives
\begin{align}
    S_1
        &= 2,
    \label{eq:cpsv16_examples_4_10_4_11_entropy_example410_entropy_one}\\
    S_2
        &\simeq 2.9544,
    \label{eq:cpsv16_examples_4_10_4_11_entropy_example410_entropy_two}\\
    S_3
        &\simeq 3.8802,
    \label{eq:cpsv16_examples_4_10_4_11_entropy_example410_entropy_three}\\
    S_4
        &\simeq 2.7839.
    \label{eq:cpsv16_examples_4_10_4_11_entropy_example410_entropy_four}
\end{align}
These values agree with the printed decimals.

The failure of saturation is exact:
\begin{align}
    I_2-I_1
        &= \frac{1}{16\ln 2}
            \ln\!\left(\frac{2^{32}7^7}{3^3 5^{20}}\right).
    \label{eq:cpsv16_examples_4_10_4_11_entropy_example410_exact_defect}
\end{align}
Because $\ln 2>0$, positivity reduces to
\begin{align}
    2^{32}7^7
        &= 3537090251849728,
    \label{eq:cpsv16_examples_4_10_4_11_entropy_example410_integer_left}\\
    3537090251849728
        &> 2574920654296875,
    \label{eq:cpsv16_examples_4_10_4_11_entropy_example410_integer_strict}\\
    2574920654296875
        &= 3^3 5^{20}.
    \label{eq:cpsv16_examples_4_10_4_11_entropy_example410_integer_right}
\end{align}
Hence $I_2>I_1$ at $p=1/4$.

The source also claims $I_2>I_1$ for every $p\ne 0,1/2$
\cite[Example~4.10]{Cirac2016MPDO_arXiv}. This exception list omits $p=1$.
At $p=1$, every spin is flipped deterministically, so the resulting state is
related to the initial Bell-pair ring by an on-site unitary and equality
returns. The certified correction is the explicit $p=1/4$ witness; the printed
universal statement must not be used as a formal hypothesis.

\section{Example~4.11: the literal periodic state}

\subsection{Exact finite-ring distribution}

The source defines diagonal two-qubit bond operators with weight matrix
\begin{align}
    W
        &= \begin{pmatrix}
            1 & 1/2\\
            1/2 & 1
        \end{pmatrix}.
    \label{eq:cpsv16_examples_4_10_4_11_entropy_example411_weight_matrix}
\end{align}
For a periodic chain of length $N$, the normalized state is the classical
distribution
\begin{align}
    \mathbb P_N(k_1,\ldots,k_N)
        &= \frac{\prod_{n=1}^{N}W_{k_n,k_{n+1}}}{Z_N},
    \label{eq:cpsv16_examples_4_10_4_11_entropy_example411_distribution}\\
    k_{N+1}
        &= k_1,
    \label{eq:cpsv16_examples_4_10_4_11_entropy_example411_periodic_boundary}\\
    Z_N
        &= \tr(W^N)
         = \frac{3^N+1}{2^N}.
    \label{eq:cpsv16_examples_4_10_4_11_entropy_example411_partition_function}
\end{align}
For $1\leq L\leq N$, summing the complementary path gives
\begin{align}
    \mu_{N,L}(k_1,\ldots,k_L)
        &= \frac{
            \left(\prod_{n=1}^{L-1}W_{k_n,k_{n+1}}\right)
            (W^{N-L+1})_{k_L,k_1}}
            {\tr(W^N)}.
    \label{eq:cpsv16_examples_4_10_4_11_entropy_example411_reduced_probability}
\end{align}
This formula includes $L=1$ and $L=N$, together with the periodic closing
factor. It is not the marginal law of an open stationary Markov chain.

Let $t(k)$ count transitions among the $L-1$ internal edges of
$k=(k_1,\ldots,k_L)$. Then
\begin{align}
    \mu_{N,L}(k)
        &= \frac{
            2^{L-t(k)-2}
            \left(3^{N-L+1}+(-1)^{t(k)}\right)}
            {3^N+1}.
    \label{eq:cpsv16_examples_4_10_4_11_entropy_example411_closed_block_probability}
\end{align}
The formal statement uses an equivalent quotient with only nonnegative
exponents. Neither block formula is printed in
Ref.~\cite{Cirac2016MPDO_arXiv}.

At $N=4$, the nonzero spectra are
\begin{align}
    \operatorname{spec}(\rho_1)
        &= \left\{(1/2)^{[2]}\right\},
    \label{eq:cpsv16_examples_4_10_4_11_entropy_example411_spectrum_one}\\
    \operatorname{spec}(\rho_2)
        &= \left\{(14/41)^{[2]},(13/82)^{[2]}\right\},
    \label{eq:cpsv16_examples_4_10_4_11_entropy_example411_spectrum_two}\\
    \operatorname{spec}(\rho_3)
        &= \left\{(10/41)^{[2]},(4/41)^{[4]},
            (5/82)^{[2]}\right\},
    \label{eq:cpsv16_examples_4_10_4_11_entropy_example411_spectrum_three}\\
    \operatorname{spec}(\rho_4)
        &= \left\{(8/41)^{[2]},(2/41)^{[12]},
            (1/82)^{[2]}\right\}.
    \label{eq:cpsv16_examples_4_10_4_11_entropy_example411_spectrum_four}
\end{align}
They give
\begin{align}
    I_2-I_1
        &= \frac{1}{82\ln 2}
            \ln\!\left(\frac{41^{82}5^{50}}
            {2^{48}13^{52}7^{112}}\right).
    \label{eq:cpsv16_examples_4_10_4_11_entropy_example411_exact_defect}
\end{align}
The exact inequality
\begin{align}
    41^{82}5^{50}
        &> 2^{48}13^{52}7^{112}.
    \label{eq:cpsv16_examples_4_10_4_11_entropy_example411_integer_inequality}
\end{align}
proves $I_2>I_1$. The difference is approximately $0.0069396647$ bits, but
the integer comparison is the proof. Thus the literal four-site periodic state
does not saturate the area law, contrary to the printed claim.

\subsection{Normalization and the likely source of the discrepancy}

Global normalization cannot repair the saturation claim. Replacing $W$ by
$cW$ for $c>0$ multiplies every unnormalized length-$N$ configuration weight
and $Z_N$ by $c^N$. The normalized state, all reduced spectra, and $I_2-I_1$
remain unchanged. In particular, replacing $W$ by $(2/3)W$ does not restore
saturation.

The source does not explain the failed finite-ring claim. One possible
explanation is confusion with the open or infinite stationary Markov chain
whose transition matrix is $(2/3)W$ and whose stationary distribution is
uniform. Define the binary entropy by
\begin{align}
    h_2(q)
        &= -q\log_2 q-(1-q)\log_2(1-q).
    \label{eq:cpsv16_examples_4_10_4_11_entropy_binary_entropy}
\end{align}
For that different process, the entropy of an open block is
\begin{align}
    H_L
        &= 1+(L-1)h_2(1/3).
    \label{eq:cpsv16_examples_4_10_4_11_entropy_open_markov_entropy}
\end{align}
This entropy is affine in $L$. The periodic closing factor in
Eq.~\ref{eq:cpsv16_examples_4_10_4_11_entropy_example411_reduced_probability}
is absent from the open process. This interpretation is suggested by the
calculation; it is not asserted in Ref.~\cite{Cirac2016MPDO_arXiv}.

\section{Formalization boundary}

Example~4.10 yields an exact corrected witness at $p=1/4$. The repeated
left-qubit label and the incomplete universal exception list must not be
formalized as printed. For Example~4.11, the printed saturation clause is
false for the literal finite periodic family. The independent claim that the
family does not have zero correlation length is verified both for the literal
diagram in Definition~4.2 of
Ref.~\cite[lines~735--739]{Cirac2016MPDO_arXiv} and for its scale-invariant
repair.

\paragraph{Local correction for the Example~4.10 tensor.}

The module \doclink{TNLean/MPS/MPDO/CPSVExample410Operator} defines the
left--right correlated flip at $p=1/4$, not the repeated-left expression. The
declaration \leanid{MPOTensor.CPSVExample410Operator.M_isMPDO} proves
positivity at every chain length. Its physical-trace transfer is the nonzero
projector
\begin{align}
    P
        &= \frac12
        \begin{pmatrix}
            1&0&0&1\\
            0&0&0&0\\
            0&0&0&0\\
            1&0&0&1
        \end{pmatrix},
    \label{eq:cpsv16_examples_4_10_4_11_entropy_example410_transfer_projector}\\
    P^2
        &= P.
    \label{eq:cpsv16_examples_4_10_4_11_entropy_example410_transfer_idempotence}
\end{align}
Consequently,
\leanid{MPOTensor.CPSVExample410Operator.M_isSourceZCL} proves source zero
correlation length for the corrected tensor.

\paragraph{Project-derived Example~4.10 spectra and entropies.}

The module \doclink{TNLean/MPS/MPDO/CPSVExample410Spectrum} identifies each
normalized reduced operator as an isometric image of a diagonal bond-window
distribution. The declarations
\leanid{MPOTensor.CPSVExample410Spectrum.charpoly_roots_one},
\leanid{MPOTensor.CPSVExample410Spectrum.charpoly_roots_two},
\leanid{MPOTensor.CPSVExample410Spectrum.charpoly_roots_three}, and
\leanid{MPOTensor.CPSVExample410Spectrum.charpoly_roots_four} prove the full
root multisets
\begin{align}
    L=1:\quad
    &\left\{0^{[0]},(1/4)^{[4]}\right\},
    \label{eq:cpsv16_examples_4_10_4_11_entropy_example410_full_roots_one}\\
    L=2:\quad
    &\left\{0^{[8]},(5/32)^{[4]},(3/32)^{[4]}\right\},
    \label{eq:cpsv16_examples_4_10_4_11_entropy_example410_full_roots_two}\\
    L=3:\quad
    &\left\{0^{[48]},(7/64)^{[4]},(3/64)^{[12]}\right\},
    \label{eq:cpsv16_examples_4_10_4_11_entropy_example410_full_roots_three}\\
    L=4:\quad
    &\left\{0^{[248]},(41/128)^{[1]},(15/128)^{[4]},
        (9/128)^{[3]}\right\}.
    \label{eq:cpsv16_examples_4_10_4_11_entropy_example410_full_roots_four}
\end{align}
These are consequences of the corrected $p=1/4$ interpretation, not formulas
printed in Ref.~\cite{Cirac2016MPDO_arXiv}.

The module \doclink{TNLean/MPS/MPDO/CPSVExample410Entropy} derives the block
entropies from these root multisets. The declarations
\leanid{MPOTensor.CPSVExample410Entropy.blockEntropy_one},
\leanid{MPOTensor.CPSVExample410Entropy.blockEntropy_two},
\leanid{MPOTensor.CPSVExample410Entropy.blockEntropy_three}, and
\leanid{MPOTensor.CPSVExample410Entropy.blockEntropy_four} identify them with
the exact correlated-flip window entropies. The declaration
\leanid{MPOTensor.CPSVExample410Entropy.mutualInfo_two_sub_one} proves
\begin{align}
    I_2-I_1
        &= \frac{1}{16}
            \ln\!\left(\frac{2^{32}7^7}{3^3 5^{20}}\right).
    \label{eq:cpsv16_examples_4_10_4_11_entropy_example410_formal_defect}
\end{align}
The declaration
\leanid{MPOTensor.CPSVExample410Entropy.mutualInfo_two_sub_one_pos} proves
positivity. The proof uses the established spectra and exact integer
comparison, not the printed decimals. Finally,
\leanid{MPOTensor.CPSVExample410Entropy.M_not_isSAL} converts the strict defect
into failure of saturation, and
\leanid{MPOTensor.CPSVExample410Entropy.M_isSourceZCL_and_not_isSAL} proves
\begin{align}
    \mathrm{IsSourceZCL}(M)
        &\mathbin{\land}\neg\mathrm{IsSAL}(M).
    \label{eq:cpsv16_examples_4_10_4_11_entropy_example410_zcl_not_sal}
\end{align}

\paragraph{Certified scalar endpoints.}

The module
\doclink{TNLean/MPS/MPDO/CPSVExamples410411Arithmetic} certifies the scalar
endpoints. The declarations
\leanid{CPSVExamples410411Arithmetic.example410_integer_inequality} and
\leanid{CPSVExamples410411Arithmetic.example411_integer_inequality} prove the
exact natural-number comparisons. The declarations
\leanid{CPSVExamples410411Arithmetic.example410_log_ratio_pos} and
\leanid{CPSVExamples410411Arithmetic.example411_log_ratio_pos} prove
positivity of the corresponding logarithmic ratios. These results entered the
project in commit \texttt{a79bfc93314f3f721abecdbc184065497672d7e5}.

\paragraph{Project-derived effective-binary finite-ring blocks.}

The module \doclink{TNLean/MPS/MPDO/CPSVExample411BinarySupport} proves the
exact internal edge product in
\leanid{MPOTensor.CPSVExample411BinarySupport.internalEdgeProduct_eq}, every
entry of $W^m$ in
\leanid{MPOTensor.CPSVExample411BinarySupport.weightMatrix_pow_apply}, and the
complementary path-sum identity in
\leanid{MPOTensor.CPSVExample411BinarySupport.sum_complementPathWeight_eq_weightMatrix_pow}.
The declaration
\leanid{MPOTensor.CPSVExample411BinarySupport.reducedBlockState_apply_eq_internal_mul_pow}
identifies the canonical reduced state with
Eq.~\ref{eq:cpsv16_examples_4_10_4_11_entropy_example411_reduced_probability}.
The declaration
\leanid{MPOTensor.CPSVExample411BinarySupport.reducedBlockState_apply} proves
the closed diagonal formula for $1\leq L\leq N$. These are finite-ring
consequences of the printed neighboring operators, not source claims. This
module remains within the effective binary model.

\paragraph{Local correction for ambient support.}

The effective binary one-site space is included in the source-facing
two-qubit site by
\begin{align}
    k
        &\longmapsto \ket{k,k}.
    \label{eq:cpsv16_examples_4_10_4_11_entropy_example411_support_inclusion}
\end{align}
It is not identified with the four-dimensional site by an equivalence. The
module \doclink{TNLean/MPS/MPDO/CPSVExample411Ambient} proves the supported
one-site slice and vanishing on unsupported local physical legs. The
declarations
\leanid{MPOTensor.CPSVExample411Ambient.mpo_ambientM_eq_singleKrausMap} and
\leanid{MPOTensor.CPSVExample411Ambient.mpo_ambientM_supported_apply} prove the
periodic-operator congruence and supported positive-length entries. The
declaration
\leanid{MPOTensor.CPSVExample411Ambient.trace_mpo_ambientM_closed_form} proves
\begin{align}
    Z_N
        &= \frac{3^N+1}{2^N}.
    \label{eq:cpsv16_examples_4_10_4_11_entropy_example411_ambient_trace}
\end{align}

\paragraph{Project-derived source-ZCL obstruction.}

The module \doclink{TNLean/MPS/MPDO/CPSVExample411SourceZCL} records
\begin{align}
    \rho^{(4)}_2(00,00)
        &= \frac{14}{41},
    \label{eq:cpsv16_examples_4_10_4_11_entropy_example411_ring_four_entry}\\
    \rho^{(3)}_2(00,00)
        &= \frac{5}{14}.
    \label{eq:cpsv16_examples_4_10_4_11_entropy_example411_ring_three_entry}
\end{align}
The declaration
\leanid{MPOTensor.CPSVExample411Ambient.reducedBlockState_supported_apply}
transports these values to the ambient two-qubit-site marginals. The
source-facing zero-correlation relation requires equality of the normalized
two-site marginals obtained from rings of lengths four and three. Since
\begin{align}
    \frac{14}{41}
        &\ne \frac{5}{14},
    \label{eq:cpsv16_examples_4_10_4_11_entropy_example411_marginal_mismatch}
\end{align}
both the effective and ambient tensors fail that relation.

Define their physical-trace transfers by
\begin{align}
    \mathcal T_M
        &:= \sum_i M^{ii},
    \label{eq:cpsv16_examples_4_10_4_11_entropy_example411_effective_transfer}\\
    \mathcal T_{\widetilde M}
        &:= \sum_i \widetilde M^{ii}.
    \label{eq:cpsv16_examples_4_10_4_11_entropy_example411_ambient_transfer}
\end{align}
The declarations
\leanid{MPOTensor.CPSVExample411BinarySupport.physTraceTransfer_M_ne_zero} and
\leanid{MPOTensor.CPSVExample411Ambient.physTraceTransfer_ambientM_ne_zero} use
\begin{align}
    \tr(\operatorname{mpo}(M,1))
        &= 2,
    \label{eq:cpsv16_examples_4_10_4_11_entropy_example411_effective_one_site_trace}\\
    \tr(\operatorname{mpo}(\widetilde M,1))
        &= 2
    \label{eq:cpsv16_examples_4_10_4_11_entropy_example411_ambient_one_site_trace}
\end{align}
to prove that both transfers are nonzero. The declarations
\leanid{MPOTensor.CPSVExample411BinarySupport.physTraceTransfer_M_not_idempotent}
and
\leanid{MPOTensor.CPSVExample411Ambient.physTraceTransfer_ambientM_not_idempotent}
then prove
\begin{align}
    \mathcal T_M^2
        &\ne \mathcal T_M,
    \label{eq:cpsv16_examples_4_10_4_11_entropy_example411_effective_nonidempotence}\\
    \mathcal T_{\widetilde M}^2
        &\ne \mathcal T_{\widetilde M}.
    \label{eq:cpsv16_examples_4_10_4_11_entropy_example411_ambient_nonidempotence}
\end{align}
Thus both fixed tensors fail the literal diagram in Definition~4.2 of
Ref.~\cite[lines~735--739]{Cirac2016MPDO_arXiv}. This argument does not use
entropy, spectra, or the saturation obstruction.

The scale-invariant statement uses a nonzero up-to-positive-scalar relation,
whereas the literal source diagram fixes the scalar to one. Both differ from
the legacy doubled-index transfer-map idempotence convention. This note makes
no claim about that convention.

\paragraph{Project-derived length-four spectra.}

The module \doclink{TNLean/MPS/MPDO/CPSVExample411Spectrum} proves the
effective nonzero root multisets
\begin{align}
    L=1:\quad
    &\left\{(1/2)^{[2]}\right\},
    \label{eq:cpsv16_examples_4_10_4_11_entropy_example411_effective_roots_one}\\
    L=2:\quad
    &\left\{(14/41)^{[2]},(13/82)^{[2]}\right\},
    \label{eq:cpsv16_examples_4_10_4_11_entropy_example411_effective_roots_two}\\
    L=3:\quad
    &\left\{(10/41)^{[2]},(4/41)^{[4]},(5/82)^{[2]}\right\},
    \label{eq:cpsv16_examples_4_10_4_11_entropy_example411_effective_roots_three}\\
    L=4:\quad
    &\left\{(8/41)^{[2]},(2/41)^{[12]},(1/82)^{[2]}\right\}.
    \label{eq:cpsv16_examples_4_10_4_11_entropy_example411_effective_roots_four}
\end{align}
The generic identity
\leanid{Matrix.singleKrausMap_charpoly_eq_X_pow_mul_of_isometry}, specialized
by
\leanid{MPOTensor.CPSVExample411Spectrum.ambient_charpoly_eq_X_pow_mul_effective},
is
\begin{align}
    \chi_{\widetilde\rho_L}(X)
        &= X^{4^L-2^L}\chi_{\rho_L}(X).
    \label{eq:cpsv16_examples_4_10_4_11_entropy_example411_charpoly_embedding}
\end{align}
The ambient block has dimension $4^L$, while the effective block has dimension
$2^L$. The added zero-root multiplicities for block lengths one through four
are therefore
\begin{align}
    2,\quad 12,\quad 56,\quad 240.
    \label{eq:cpsv16_examples_4_10_4_11_entropy_example411_zero_multiplicities}
\end{align}
These finite-periodic spectra are project-derived, not source statements.

\paragraph{Project-derived entropy and saturation obstruction.}

The module \doclink{TNLean/MPS/MPDO/CPSVExample411Entropy} proves that the
effective binary tensor is an MPDO and evaluates the four block entropies from
the root multisets. In natural-logarithm units,
\begin{align}
    I_2-I_1
        &= \frac{1}{82}
            \ln\!\left(\frac{41^{82}5^{50}}
            {2^{48}13^{52}7^{112}}\right),
    \label{eq:cpsv16_examples_4_10_4_11_entropy_example411_formal_defect}\\
    I_2-I_1
        &> 0.
    \label{eq:cpsv16_examples_4_10_4_11_entropy_example411_formal_defect_positive}
\end{align}
Dividing by $\ln 2$ gives the bit-valued formula. The strict inequality reuses
\leanid{CPSVExamples410411Arithmetic.example411_log_ratio_pos}; it does not
recompute the spectra or integer comparison. The effective binary tensor
therefore fails SAL. Contraposition of SAL descent through the physical
isometry gives the same conclusion for the ambient tensor
\leanid{MPOTensor.CPSVExample411Ambient.ambientM}.

\paragraph{Remaining scope restrictions.}

For Example~4.11, the ambient operator congruence, trace equality, exact
four-site spectra, entropy identities, finite-ring saturation obstruction, and
source-facing zero-correlation-length obstruction are complete. The printed
non-zero-correlation-length clause is verified for the literal finite periodic
family, while the printed saturation clause is false. No conclusion is made
about thermodynamic behavior or the legacy doubled-index convention.

For Example~4.10, the corrected fixed-$p=1/4$ tensor, its MPDO property,
source zero correlation length, complete one- through four-site spectra, block
entropies, positive mutual-information defect, and non-SAL conclusion are
complete. These results do not establish the printed universal statement in
$p$, and $p=1$ is an additional exception.

Example~4.12 lies outside this entropy correction. Its statement,
dependencies, and separate normalization question remain unchanged.

\section{Verdict}

For Example~4.10, the repeated left-qubit label requires a local correction to
the left--right correlated flip. At $p=1/4$, the corrected tensor has source
zero correlation length, and its complete one- through four-site spectra give
the exact block entropies. The positive defect $I_2-I_1$ proves failure of
area-law saturation. These spectra and exact entropy identities are
project-derived; the source prints only the channel and decimal entropies. The
printed universal exception list must also include $p=1$.

For Example~4.11, the printed separation is only partly correct for the literal
periodic state. The exact finite-marginal mismatch verifies failure of
source-facing zero correlation length. The four-site result $I_2>I_1$ refutes
the printed saturation claim. Global normalization cannot repair this failure.
The open or infinite Markov-chain interpretation remains a suggested
explanation, not a source statement.

\bibliographystyle{alpha}
\bibliography{references}

\end{document}
